\chapter{Supplementary Results}
\label{app:supplementary_results}

\section{Rust Anchor Repetition Summary}

The repeated cold-cache anchor set was used to stabilise the selector redesign before the full matrix was rerun. Table~\ref{tab:anchor-summary} summarises the mean-best backend for the three anchor models under each format.

\begin{table}[h]
\caption{Mean-best explicit backend on the five-run cold-cache anchor set.}
\label{tab:anchor-summary}
\centering
\begin{tabular}{lll}
\toprule
Format & Model & Mean-best backend \\
\midrule
SafeTensors & Qwen3-0.6B & sync \\
SafeTensors & Qwen3-8B & io\_uring \\
SafeTensors & Qwen3-32B & io\_uring \\
ServerlessLLM & Qwen3-0.6B & async/io\_uring crossover \\
ServerlessLLM & Qwen3-8B & io\_uring \\
ServerlessLLM & Qwen3-32B & io\_uring \\
\bottomrule
\end{tabular}
\end{table}

This anchor set was useful because it exposed the regime split that the earlier single-rule selector had missed.

\section{vLLM Decode Supplement for Qwen3-8B}

Table~\ref{tab:vllm-extra} records the steady-state decode results for Qwen3-8B under warm and cold cache conditions. Unlike the load-only and TTFT experiments, the spread here is very small, reinforcing the claim that backend choice matters much more for model initialisation than for steady-state decode.

\begin{table}[h]
\caption{Qwen3-8B steady-state decode latency (ms) by loader and cache condition.}
\label{tab:vllm-extra}
\centering
\begin{tabular}{lcc}
\toprule
Loader & Warm decode & Cold decode \\
\midrule
native & 58.18 & 58.18 \\
ts\_safetensors\_default & 58.12 & 58.33 \\
ts\_safetensors\_sync & 58.16 & 58.23 \\
ts\_safetensors\_io\_uring & 58.11 & 58.44 \\
ts\_serverlessllm\_default & 58.06 & 58.14 \\
ts\_serverlessllm\_sync & 58.39 & 58.13 \\
ts\_serverlessllm\_io\_uring & 58.33 & 58.10 \\
\bottomrule
\end{tabular}
\end{table}

The spread between the best and worst cold decode result is only 0.34~ms, which is negligible relative to the much larger differences observed in model initialisation.

\section{Large-Model vLLM Rerun Note}

The initial vLLM matrix produced failures for the 14B and 32B models. These failures were caused by conservative harness settings rather than by checkpoint backend limitations. After increasing \texttt{gpu\_memory\_utilization} and lowering \texttt{max\_model\_len}, all failed rows were rerun successfully, yielding a completed matrix with 672 successful rows and no remaining failures.

This rerun matters because it shows that large-model feasibility in vLLM depends on serving-engine memory policy as well as loader speed. In other words, the benchmark harness itself can become a limiting factor even when the underlying checkpoint loading strategy is sound.
